% II. Анализ на възможните решения и обосновка на избора.
% Всеки избор тук трябва да има отхвърлена алтернатива и записана причина за отхвърлянето.
\section{Анализ на възможните решения и обосновка на избора}
\label{sec:alternatives}

\subsection{Семейство на протокола}

Първият избор е между двете семейства протоколи за вътрешна маршрутизация. Протокол по
вектор на разстоянието би бил по-кратък за реализация, но при него всеки възел вижда
само обобщени разстояния и няма представа за топологията. Измерването на преходни цикли
тогава изисква външен наблюдател, който събира таблиците на всички възли и възстановява
глобалната картина. При протокол по състояние на връзките всеки възел вече притежава
описание на топологията, така че разминаването между представите на два възела е
непосредствено наблюдаема величина. Избран е подходът по състояние на връзките именно
защото прави измерваната величина видима отвътре.

\subsection{Модел на времето}

Вторият избор е между три модела на времето. Работа изцяло върху реални интерфейси дава
най-висока достоверност, но не дава повторяемост: моментът на отказа не се възпроизвежда
и разсейването на резултатите скрива ефекта, който се изследва. Работа изцяло в
симулатор дава пълна повторяемост, но проверява модела, а не програмата, която после ще
работи в мрежа. Приет е трети вариант: протоколната логика не получава достъп до
часовник и до сокети, а взаимодейства с абстракция, зад която стои или дискретно
събитийна опашка с виртуален часовник, или системният часовник и реалните сокети.
Едни и същи изходни файлове на протокола се компилират в двата режима. Цената на този
избор е допълнителна абстракция и невъзможност да се използва пряко системно извикване
в протоколния код. Ползата е, че измереното поведение принадлежи на реализацията, а не
на нейн симулационен двойник.

\subsection{Език и инструментариум за изграждане}

Разгледани са C++, Rust и Go. Go предлага най-кратък път до работеща мрежова програма,
но неговата среда за изпълнение управлява планирането и събирането на паметта, което
внася закъснения, невъзможни за отстраняване от измерването. Rust дава силни гаранции за
безопасност на паметта, но еталонните реализации, спрямо които се сверява коректността,
са на C и интеграцията с тях изисква допълнителен слой. Избран е C++ по стандарта от
2020 година~\cite{iso_cpp20}: детерминирано управление на ресурсите, пряка съвместимост с
интерфейсите на операционната система и наличие на конструкции, които правят
абстракцията върху времето безплатна по време на изпълнение. За изграждане е избран
CMake, защото трябва да произвежда еднакъв резултат върху Windows, Linux и macOS.

\subsection{Еталон за сверяване на коректността}

Като еталон са разгледани FRRouting и BIRD. И двата са зрели реализации на OSPF, които
работят срещу оборудване на производители и следователно тълкуват~\cite{rfc2328} по
начин, приет от практиката. Еталонът би се използвал за проверка на коректността, тоест за
съвместимост на форматите на пакетите и на съдържанието на изчислените маршрути, а не за
сравнение по бързина. Сравнение по бързина би било подвеждащо, защото обхватът на
реализацията в този проект е подмножество от този на еталона.

Този избор беше преразгледан по време на реализацията и 	extbf{отхвърлен}. Съвместимостта
с външна реализация изисква форматът по мрежата да бъде този на OSPFv2 по~\cite{rfc2328},
включително двойките пакети за обзор на базата с водещ и подчинен, областите и типовете
обявления, които настоящият обхват не използва. Реализацията на кодирането е съизмерима по
обем с останалата част от проекта и не допринася за измерваната величина, защото времето за
сходимост не зависи от подредбата на полетата в пакета. Вместо външен еталон е използван
вътрешен: маршрутните таблици на всички възли се сравняват с таблици, изчислени независимо
върху истинската топология. Разликата в силата на двете проверки е отчетена в
раздел~
ef{sec:results}, а не премълчана.

\subsection{Обхват спрямо учебното съдържание}

Дисциплината обхваща и безжични, клетъчни, спътникови и ATM мрежи. Те са включени в
раздел~\ref{sec:literature} като контекст, но не в реализацията. Моделирането на канал
със загуби и променливо закъснение изисква обосноваване на самия модел и измерената
сходимост тогава зависи от достоверността на този модел. Проектът избира по-тесен обхват
с ясно измерима величина пред по-широк обхват с недоказуеми числа.
